% ======================================================================
% SECTION 4: DISCUSSION
% Five thematic subsections: significance of the 60-second target,
% FDA framing, on-prem LLM safety advantages, head-start framing for
% the industry, and practical insights on how close the application
% is to real life today.
% Location: 2030-gbm-1min/paper/full-paper/sections/discussion.tex
% ======================================================================

\subsection{Significance of the 60-Second Computational Target}
\label{sec:discussion-significance}

The 60-second resection target is a computational checkpoint, not a
clinical claim. The realized quality benefit is a 240x to 480x
compression of operative duration relative to the published 4 to 8
hour single-arm baseline for an IDH-wildtype glioblastoma resection
\cite{Stupp2005Glioblastoma, Sanai2011GBMResection}. Under the frozen
composite score formula (quality 0.40, time 0.25, cost 0.20,
safety 0.10, patient experience 0.05) the compression translates
into measurable patient-experience and safety gains as summarized in
Table \ref{tab:significance}.

\begin{table}[H]
\centering
\caption{Realized quality and safety benefits relative to the 4 to 8
hour baseline. Source files are at
2030-gbm-1min/outputs/comparison/ and
outputs/comparison_robot_vs_human/.}
\label{tab:significance}
\begin{tabular}{>{\raggedright\arraybackslash}p{2.8cm}
                >{\raggedright\arraybackslash}p{3.3cm}
                >{\raggedright\arraybackslash}p{2.7cm}
                >{\raggedright\arraybackslash}p{2.4cm}
                >{\raggedright\arraybackslash}p{2.4cm}}
\toprule
\textbf{Dimension} & \textbf{Baseline 4-8 h} & \textbf{60 s target} & \textbf{Rel. improvement} & \textbf{Source} \\
\midrule
Procedure duration & 14,400 to 28,800 s & 60 s & 240x to 480x & comparison_report.md \\
Time score (0 to 1) & 0.20 & 1.00 & 5x & metrics/ \\
Per-arm tip force max & 15.0 N (single arm) & 5.0 N (per arm) & 3x tighter & summary.json \\
Cumulative cross-arm force & not specified & 12.0 N envelope & new floor & arm_heartbeat.cpp \\
E-stop latency & 50 ms & 5 ms & 10x & robot_loop_4arm.cpp \\
Composite (robot mean) & 70.35 (human) & 88.53 (robot) & +25.8 percent & comparison.json \\
\bottomrule
\end{tabular}
\end{table}

The per-arm contribution analysis at
outputs/viz/per_arm_contribution.png shows arm 1 dominates resection
volume (32,400 mm cubed mean) while arms 2 to 4 stay within the 30
percent target balance band on coagulation seconds, suction volume,
and imaging frames respectively. This balance is the mechanical
signature of the on-premises LLM single-robot-error-minimization
thesis and is the property that a 1-arm baseline cannot produce
\cite{rosa-one-brain, Lefranc2014ROSAStereotaxy, kawchak_2026_19994945,
repo-robotic-surgeries}.

\subsection{Respectful Yet Opportunistic FDA Framing}
\label{sec:discussion-fda}

The 60-second glioblastoma resection is positioned as an
opportunistic extension of the FDA 28 April 2026 Real-Time Clinical
Trials (RTCT) proof-of-concept program \cite{fda2026realtime} from
pharmacology into the surgical theater. The two named RTCT pilots
(AstraZeneca's TRAVERSE in mantle cell lymphoma and Amgen's
STREAM-SCLC in small cell lung carcinoma) cover pharmacology only;
the surgical theater is the next opportunistic surface to extend.
The United States needs to remain Number 1 in the world regarding
patient safety, efficacy, and speed benefits in oncological robotic
surgeries in clinical trials, and the on-premises LLM thesis is one
route to that goal. The work is framed as an extension, not a
critique.

The appropriate regulatory pathway for an on-premises LLM that emits
control commands is the FDA Software as a Medical Device (SaMD)
framework \cite{fda-samd}. The standards floor is set by IEC 80601-
2-77 (per-arm 5.0 N tip and 1.0 N lateral force limits, 5 ms E-stop
budget) \cite{iec-80601-2-77} and IEC 62304 (safety-critical
software lifecycle) \cite{iec-62304}. The six 4-entity LLM
tournament rounds at outputs/comparison/ and the six mixed
tournament rounds at outputs/comparison_robot_vs_human/ preserve the
structural-time-dimension caveat in every rationale. The caveat is
the operational signature of a respectful framing: the 1-minute
robot trivially beats the 1-hour human-baseline scenario on the
time dimension, but the comparison is flagged as structural and
not a fair pairwise comparison.

\subsection{On-Premises LLM Single-Robot Error Minimization}
\label{sec:discussion-llm}

The mechanism that minimizes single-robot error potential is
cross-arm safety on the deterministic 1 kHz heartbeat broadcast bus
implemented at
kevinkawchak/robotic-surgeries/2030-gbm-1min/src/coordination/arm_heartbeat.cpp.
The 32-byte heartbeat frames carry per-arm state at the 1 ms
deadline; the 3 ms watchdog parks any arm that misses three
consecutive frames; the cumulative tip force is capped at 12 N on
the patient frame; per-arm tip force is capped at 5.0 N and per-arm
lateral force at 1.0 N \cite{iec-80601-2-77, iec-62304}. The 5 ms
E-stop budget plus 100 microsecond emergency arm-park trigger
enforces the safety envelope at the actuator interface.
Table \ref{tab:cross-check} renders the 4-arm cross-check matrix.

\begin{table}[H]
\centering
\caption{4-arm cross-check matrix. Each arm cross-checks its three
peers on the deterministic 1 kHz heartbeat bus; a violation triggers
the 5 ms E-stop with the 100 microsecond emergency arm-park.}
\label{tab:cross-check}
\begin{tabular}{>{\raggedright\arraybackslash}p{1.0cm}
                >{\raggedright\arraybackslash}p{2.6cm}
                >{\raggedright\arraybackslash}p{3.2cm}
                >{\raggedright\arraybackslash}p{3.4cm}
                >{\raggedright\arraybackslash}p{3.0cm}}
\toprule
\textbf{Arm} & \textbf{Tool} & \textbf{Primary sensors} & \textbf{Cross-check on peers} & \textbf{Fault response} \\
\midrule
1 & hybrid u-w-p & tip force, removal rate & cumulative cap, arm 4 imaging & throttle, then park \\
2 & bipolar coag plus irrigation & bipolar current, irrigation flow & arm 1 hemostasis demand & coag halt, then park \\
3 & suction plus collection & flow rate, occlusion & arm 1 debris yield & flush, then park \\
4 & 0.5 T iMRI plus 5-ALA plus US & imaging frames, fluorescence & arm 1 cut margin, arm 2 coag site & re-image, then park \\
\bottomrule
\end{tabular}
\end{table}

The on-premises LLM control loop ASCII diagram at
2030-gbm-1min/outputs/diagrams/ instantiates the thesis: an LLM that
holds the entire repository (12 instruction files, every kinematics
and schema specification, all 16 iteration aggregates) in a single
1M token context window \cite{claude-opus-47, claude-code} emits per-
arm xyz commands at the 1 kHz tick. Commands that would violate the
cumulative tip-force envelope are rejected at the heartbeat boundary
before they reach the actuator. On-premises inference keeps patient
health information on site (no cloud transit) when the deployment
uses Ollama \cite{ollama} or vLLM \cite{vllm} as the backend; the
Anthropic API remains a viable default for non-PHI development.

\subsection{Head Start: Checking Steps the Industry No Longer Has to Take}
\label{sec:discussion-headstart}

The 16-iteration deterministic sweep, the 4-entity LLM tournament,
the cross-platform build recipes, and the 60-second simulation
pipeline are checking steps the industry no longer has to take to
evaluate the feasibility of an on-premises LLM controlled 4-arm
parallel stereotactic neurosurgical robot. Table \ref{tab:headstart}
enumerates the checking steps and the corresponding project status,
sourced from the narrative reports at
2030-gbm-1min/outputs/reports/run_summary.md,
outputs/reports/process_log.md, and outputs/reports/final_report.md.

\begin{table}[H]
\centering
\caption{Checking steps already taken under this project. Industry
status reflects the published state of the art at time of writing.}
\label{tab:headstart}
\begin{tabular}{>{\raggedright\arraybackslash}p{3.0cm}
                >{\raggedright\arraybackslash}p{2.6cm}
                >{\raggedright\arraybackslash}p{3.0cm}
                >{\raggedright\arraybackslash}p{3.2cm}
                >{\raggedright\arraybackslash}p{2.4cm}}
\toprule
\textbf{Checking step} & \textbf{Industry status} & \textbf{Project status} & \textbf{Source artifact} & \textbf{Downstream effect} \\
\midrule
7-DOF DH parameters per arm & not published for 1-min target & complete & config/kinematics_4arm.yaml & xyz mapping ready \\
1 kHz heartbeat broadcast bus & not standardized & complete & src/coordination/arm_heartbeat.cpp & cross-arm safety \\
Per-arm force envelope and cross-check matrix & per-arm only & cumulative cap plus cross-check & arm_heartbeat.cpp plus iterations/ & 12 N safety \\
16-iteration deterministic sweep with bit determinism & no published sweep at this scale & complete & outputs/iterations/ & reproducible \\
Structured LLM tournament with on-prem inference & cloud-only & complete & outputs/comparison/ + outputs/comparison_robot_vs_human/ & on-prem ready \\
\bottomrule
\end{tabular}
\end{table}

The progression from author work on patient matching
\cite{kawchak_2025_17774560}, ten-year meta-analysis at scale
\cite{kawchak_2025_15549831}, drug synergy machine learning
\cite{kawchak_2025_17614396}, and 1-hour 4-sponsor simulations
\cite{kawchak_2026_19994945} into the 60-second 4-arm robotic target
\cite{repo-robotic-surgeries, one-minute-variant-instructions}
materializes the head-start framing as a chain of checking steps
already executed.

\subsection{Practical Insights: How Close to Real Life Today}
\label{sec:discussion-practical}

The 60-second glioblastoma robotic resection is not feasible with
current commercial robotics. The simulation pipeline is nevertheless
reproducible bit for bit from seed 20260510 and runs on commodity
hardware (M3 Ultra wall-clock 26 to 32 s per iteration, A100 12 s
per iteration; read 2030-gbm-1min/README.md). The gaps to real life
are:

(a) The hypothetical NeuroSpeed 1.0 hardware does not exist. The
robot needs 5x to 200x improvements over Medtronic ROSA ONE Brain
v3.0 across velocity, acceleration, E-stop latency, positioning
accuracy, and force resolution \cite{rosa-one-brain}.

(b) The synthetic patient PAT-GBM-0001 is not real. There is no
institutional review board oversight, no informed consent, and no
real electronic health record signal behind any artifact in the
2030-gbm-1min/ tree.

(c) The deterministic 16-iteration sweep has not been validated
against any real-patient cohort or the TRIPOD+AI / CREMLS reporting
floor \cite{Collins2024TRIPODAI, ElEmam2024CREMLS}.

The gaps that are close to closing today are:

(i) The on-premises LLM inference backends are production-grade for
the 1M token context window \cite{ollama, vllm, claude-opus-47}.

(ii) The Apache Parquet plus DuckDB analytics stack is production-
grade for per-iteration aggregates and cross-iteration analytics
\cite{apache-arrow, duckdb}.

(iii) The Zenodo L0 deposition protocol is production-grade at DOI
10.5281/zenodo.18445179 \cite{zenodo, repo-robotic-surgeries}.

(iv) The repository CI gates (ruff format, ruff check, yamllint) on
Python 3.10, 3.11, and 3.12 already cover the Python plus YAML plus
Markdown surface of the simulation pipeline.

The robotic-surgeries simulation pipeline is therefore a "ready-to-
port" software substrate: the moment a real NeuroSpeed-class robot
is available, the on-premises LLM control loop can be re-pointed at
it with minimal additional engineering. The forward path is detailed
in \S\ref{sec:limits-future}.
